\documentclass[conference]{IEEEtran}
\IEEEoverridecommandlockouts
\usepackage{cite}
\usepackage{amsmath,amssymb,amsfonts}
\usepackage{graphicx}
\usepackage{textcomp}
\usepackage{xcolor}
\usepackage{url}

\def\BibTeX{{\rm B\kern-.05em{\sc i\kern-.025em b}\kern-.08em
    T\kern-.1667em\lower.7ex\hbox{E}\kern-.125emX}}

\begin{document}
\title{Perancangan IP I2S Transmitter AXI4-Lite pada SoC MicroBlaze V menggunakan Verilog}

\author{\IEEEauthorblockN{Muhammad Muqtada Alhaddad}
\IEEEauthorblockA{Program Studi Teknologi Informasi\\
Universitas Gadjah Mada\\
Yogyakarta, Indonesia\\
muqtada@example.com}
}

\maketitle

\begin{abstract}
This paper presents the implementation of a custom I2S transmitter IP core in Verilog that is controlled through an AXI4-Lite register interface on a MicroBlaze V soft-core SoC. The target platform is a Basys3 FPGA board driving an external PCM5102A DAC. Audio timing is generated from a 100 MHz system clock using a numerically controlled oscillator, while the firmware streams stereo samples through memory-mapped register writes and status polling. The final design supports Philips I2S framing, 32-bit slots with 24-bit audio, and programmable integer sampling rates without relying on FIFO, DMA, or a licensed audio IP block. Simulation and hardware validation confirm correct waveform generation, stable audio output, and a compact implementation suitable as reusable soft IP.
\end{abstract}

\begin{IEEEkeywords}
I2S, AXI4-Lite, MicroBlaze V, FPGA, Verilog, PCM5102A, soft IP
\end{IEEEkeywords}

\section{Introduction}
Digital audio on FPGA is attractive when deterministic timing, configurable interfaces, and software-controlled peripheral access are required. The I2S protocol remains a practical transport for PCM audio because it is simple, standardized, and directly compatible with external DACs\cite{philips_i2s_1986}. For SoC control, AXI4-Lite provides a compact memory-mapped interface for custom peripherals\cite{arm_amba_axi_2010}.

This work follows the thesis implementation path for a custom I2S transmitter IP that is integrated into a MicroBlaze V based SoC on Basys3 and connected to a PCM5102A DAC\cite{xilinx_microblaze_ref_2021,xilinx_basys3_ref,pcm5102a_datasheet}. The goal is not only to produce correct I2S waveforms, but also to keep the peripheral reusable as a soft IP block in Vivado designs\cite{xilinx_vivado_design_suite_2021}.

The main contribution is a complete hardware-software flow: a Verilog I2S transmitter core, an AXI4-Lite register map, an internal phase-accumulator timing generator, and bare-metal firmware that streams audio samples using status polling. The implementation stays transparent at RTL level and avoids dependence on a licensed vendor audio block.

\section{Background and Related Concepts}
\subsection{I2S audio framing}
I2S uses three signals: bit clock, word select, and serial data. Philips I2S requires a one-bit delay between the word-select transition and the first data bit, and the design in this paper preserves that rule\cite{philips_i2s_1986}. The transmitter operates in stereo mode with 32-bit slots, while the active sample width is 24 bits.

\subsection{AXI4-Lite register access}
AXI4-Lite is used as a simple register-access protocol for control and status traffic\cite{arm_amba_axi_2010,xilinx_axi_ref_2017}. In this design, the software writes audio samples and control settings to memory-mapped registers, then waits for a status transition that indicates the hardware has latched the sample pair.

\subsection{Clock generation}
The audio timing is derived from a 100 MHz system clock using a numerically controlled oscillator. NCO-style phase accumulation is a standard way to generate programmable output frequencies in digital systems\cite{vankka_dds_2001,ao2012nco}. Here, that technique is used to derive BCLK and MCLK without requiring a separate audio PLL.

\section{System Design}
\subsection{Hardware architecture}
The final system contains four main parts: MicroBlaze V, AXI interconnect, the custom I2S IP, and the external PCM5102A DAC. The processor writes to the I2S peripheral through AXI4-Lite, and the peripheral outputs \texttt{i2s\_mclk}, \texttt{i2s\_bclk}, \texttt{i2s\_ws}, and \texttt{i2s\_data}. UART is reserved for debug output during development.

\begin{figure}[t]
\centering
\IfFileExists{../thesi2s/contents/figures/processor-uart-timer-axi-i2s-architecture-whole-diagram.png}{%
\includegraphics[width=0.92\columnwidth]{../thesi2s/contents/figures/processor-uart-timer-axi-i2s-architecture-whole-diagram.png}%
}{}
\caption{System architecture linking MicroBlaze V, UART, AXI4-Lite, and the custom I2S transmitter IP.}
\label{fig:system-architecture}
\end{figure}

\subsection{Register map}
Table~\ref{tab:register-map} summarizes the peripheral interface.

\begin{table}[t]
\caption{AXI4-Lite register map of the I2S transmitter}
\label{tab:register-map}
\centering
\begin{tabular}{|l|c|p{4.1cm}|}
\hline
\textbf{Register} & \textbf{Offset} & \textbf{Function} \\
\hline
DATA\_LEFT  & 0x00 & Left-channel audio sample written by software \\
\hline
DATA\_RIGHT & 0x04 & Right-channel audio sample written by software \\
\hline
CONTROL     & 0x08 & Enable, mute, and programmable sampling-rate control \\
\hline
STATUS      & 0x0C & Frame-latch status bit used for software synchronization \\
\hline
\end{tabular}
\end{table}

The control path is intentionally small so the firmware can interact with the peripheral using ordinary load/store operations rather than a driver stack. This keeps the software side suitable for bare-metal Vitis projects.

\subsection{Clock and frame timing}
The implementation uses a 100 MHz reference clock and generates the audio clocks from phase accumulators. The design aligns the channel boundary with the 64-bit stereo frame, meaning one frame contains 32 BCLK periods for the left channel and 32 for the right channel. A status toggle marks the moment when the hardware has moved the current sample pair into its internal shadow registers.

\section{Implementation}
\subsection{RTL structure}
The RTL is organized as a soft IP core written in Verilog\cite{palnitkar_verilog_2003}. The core contains three functional blocks: AXI4-Lite handshake logic, a register file, and an I2S serializer. The register file stores the audio samples and control bits, while the serializer emits the Philips I2S bit stream. The clock generator and frame counter are kept inside the core so the design is self-contained.

\subsection{Firmware flow}
The bare-metal firmware initializes the peripheral, writes the control word, waits for the frame-status transition, and then pushes the next stereo sample pair. The software uses status polling instead of delays or interrupts so that sample submission stays synchronized with the frame boundary.

This approach avoids the drift that would occur if the software tried to approximate audio timing using sleep calls alone. It also keeps the transfer path deterministic across sample-rate changes.

\subsection{Supported operating modes}
The implementation supports the following modes:
\begin{itemize}
\item Philips I2S framing with one-bit delay.
\item Stereo output with 32-bit slots.
\item 24-bit active sample width.
\item Programmable integer sampling rates, including standard audio rates such as 44.1 kHz, 48 kHz, 88.2 kHz, and 96 kHz.
\item Mute and enable control through the AXI4-Lite control register.
\end{itemize}

\section{Results and Discussion}
\subsection{Functional validation}
The RTL was verified in simulation before board deployment. The checks covered AXI4-Lite read/write behavior, BCLK and WS generation, one-bit I2S delay, stereo frame alignment, and mute behavior. After the RTL passed simulation, the bitstream was implemented on Basys3 and the DAC output was observed on the physical board.

Figure~\ref{fig:system-architecture} summarizes the full hardware-software chain used in the implementation.

The hardware test confirmed that the PCM5102A produced clean audio output for the intended modes. In particular, the final firmware avoided the sign-bit corruption and timing drift issues that were encountered during development.

\subsection{Synthesis summary}
Table~\ref{tab:synthesis} summarizes the reported implementation results.

\begin{figure}[t]
\centering
\IfFileExists{../thesi2s/contents/figures/utilization-summary.png}{%
\includegraphics[width=0.92\columnwidth]{../thesi2s/contents/figures/utilization-summary.png}%
}{}
\caption{Resource utilization summary of the implemented SoC and I2S transmitter.}
\label{fig:utilization-summary}
\end{figure}

\begin{table}[t]
\caption{Reported implementation summary}
\label{tab:synthesis}
\centering
\begin{tabular}{|l|c|}
\hline
\textbf{Metric} & \textbf{Value} \\
\hline
Slice registers and logic & 903 Slice \\
\hline
LUT usage & 2435 LUT \\
\hline
Flip-flop usage & 2330 FF \\
\hline
BRAM usage & 32 BRAM \\
\hline
Total power & 0.209 W \\
\hline
WNS at 100 MHz & 0.064 ns \\
\hline
Estimated Fmax & 100.64 MHz \\
\hline
\end{tabular}
\end{table}

These numbers indicate that the design is lightweight enough for Artix-7 class devices while still providing a reusable audio peripheral and software control path.

Figure~\ref{fig:utilization-summary} provides a compact view of the implementation footprint reported by the thesis build.

\subsection{Discussion}
The main advantage of the final architecture is that the peripheral does not depend on a third-party audio IP license. The core remains transparent at RTL level, and the register map is simple enough to be driven from bare-metal firmware. Compared with a FIFO- or DMA-based design, this implementation is easier to understand and sufficient for modest audio streaming workloads.

The tradeoff is that software must respect the frame boundary and keep up with the status polling loop. That is acceptable for the intended use case, but a higher-throughput system would benefit from FIFO buffering or DMA support.

\section{Conclusion}
A custom I2S transmitter IP was successfully designed in Verilog and integrated into a MicroBlaze V SoC using AXI4-Lite control. The design generates Philips I2S timing, supports stereo 24-bit audio in 32-bit slots, and uses a phase-accumulator clocking scheme to support programmable integer sampling rates. Hardware validation with a PCM5102A DAC and synthesis results on Basys3 show that the implementation is compact, reusable, and practical for FPGA-based audio experiments.

Future work can extend the core with TX FIFO buffering, interrupt-driven transfers, or DMA support for higher sustained throughput.

\bibliographystyle{IEEEtran}
\bibliography{references}

\end{document}
